\scanpage{20}\mspage{19}
% slanted
\margin{partie de \deleted{\unsure{Th}} \unsure{th.}}
\deleted{\uline{\unsure{Prop}}}\inserted{\uline{Prop} Soit $X,Y$} Deux parties fermées modérées, \unsure{en} pos. rel. modérée. Alors $X\cap Y$ est une partie \inserted{fermée} modérée, et
\[
\boxed{(X\cap Y)^{\unsure{\cdot|\prime}}\subset X^{\cdot}\cup Y^{\cdot}}
\]
De plus, \inserted{$X$ et $Y$ ont une borne sup. \unsure{dans|en}} \deleted{\unsure{parmi} \unsure{toutes les} parties} l'ens des parties modérées ordonnées par \deleted{fermées modérées qui \unsure{contient} $X$} l'inclusion.
% slanted
\margin{\deleted{\unsure{Dém}} \uline{Dém} double récurrence sur nb de comp. conn. de $X$ et de $Y$ (pour l'existence du \unsure{Sup|sup} et la \unsure{question des frontières})}
\deleted{$X$ et $Y$, et on} On écrit cette partie $X\underset{\mathrm{mod}}{\cup}Y$, On a
\[
\boxed{(X\underset{\mathrm{mod}}{\cup}Y)^{\cdot}\subset X^{\cdot}\cup Y^{\cdot}}
\]
% slanted
% Two nearly parallel strokes separate the left notes, immediately above Rem.
\margin{Rem, \unsure{tte} partie modérée $Z$ en pos. rel. \unsure{mod.} avec $X$ et $Y$, l'est aussi avec $X\cap Y$ et $X\underset{\mathrm{mod}}{\cup}Y$}

\uline{Cor} Soit $(X_i)_{i\in I}$ une famille finie de \unsure{partie|parties} fermées modérées mutuellement en pos. \deleted{rel} modérée. Alors $\bigcap_i X_i$ est une partie modérée, et les $X_i$ ont une borne sup. \unsure{ds|dans} l'ens des parties fermées modérées. On a
% slanted
\margin{\uline{Attention} pour une partie \uline{\unsure{quelc|seule}} $A$ de $\unsure{L|\mathscr{L}}$, il n'y a \uline{pas} en général d'\unsure{adh.} modérée fermée, i.e. de ``plus petit'' ens \unsure{f.} modéré $X$ tel que $X\supset A$. C'est déjà faux ds le cas standart $L=[0,1]$ ! (Prendre $A=\{\unsure{0|1},1/2,1/3\ldots\}$}
\[
\left(\bigcup_{\substack{\mathrm{mod}\\ i}} X_i\right)^{\cdot}\subset\bigcup_i X_i^{\cdot}
\]

\uline{Ex}
\[
[a,b]\underset{\mathrm{mod}}{\cup}[b,c]=[a,c]
\]
mais en général
\[
[a,b]\cup[b,c]\ne[a,c]
\]
(\unsure{prendre|quand} $a$ plus \unsure{c|e} petit élément, $c$ plus grand élément \deleted{\unsure{et}} \unsure{qqconq} \deleted{\unsure{Et en ce} $\unsure{=}$}, $b$ \deleted{\unsure{est}} comp. à tout \inserted{élément de $L$.} Si c'est
